% Generated draft: requires mathematical and source review.
\documentclass[11pt]{article}
\usepackage[margin=25mm]{geometry}
\usepackage{fontspec}
\setmainfont{Latin Modern Roman}
\usepackage{amsmath,amssymb,mathtools,unicode-math}
\setmathfont{Latin Modern Math}
\usepackage{enumitem,needspace,xurl,hyperref}
\setlength{\parindent}{0pt}
\setlength{\parskip}{5pt}
\setlength{\emergencystretch}{3em}
\newcommand{\N}{\mathbb{N}}
\newcommand{\Z}{\mathbb{Z}}
\newcommand{\Q}{\mathbb{Q}}
\newcommand{\R}{\mathbb{R}}
\newcommand{\C}{\mathbb{C}}
\newcommand{\F}{\mathbb{F}}
\newcommand{\A}{\mathbb{A}}
\newcommand{\PP}{\mathbb{P}}
\newcommand{\E}{\mathbb{E}}
\newcommand{\eps}{\varepsilon}
\newcommand{\dd}{\,\mathrm{d}}
\begin{document}
\section*{UnsolvedMath ID: 30003783}
\textbf{Source problem code:} catalogue entry
\subsection{Definitions and mathematical statement}
The supplied question is: ``Does a rigidity theorem hold for quasi-isometric embeddings, or even coarse embeddings, between corresponding warped-cone level sequences?'' The record does not specify the actions, the correspondence between levels, or the conclusion meant by rigidity. Consequently, it does not determine a single fully quantified conjecture.

The following definitions describe the terminology conditionally; they do not supply missing hypotheses for the intended problem. In a customary warped-cone construction, let a finitely generated group \(\Gamma\), with finite symmetric generating set \(S\), act by homeomorphisms on a compact metric space \((X,d_X)\). For \(t>0\), the warped metric \(d_{\Gamma,t}\) is the largest metric on \(X\) satisfying
\[
 d_{\Gamma,t}(x,x')\leq t\,d_X(x,x'),\qquad
 d_{\Gamma,t}(x,sx)\leq 1
 \quad(x,x'\in X,\ s\in S).
\]
The metric space \((X,d_{\Gamma,t})\) is a level of the warped cone. A second action \(\Lambda\curvearrowright(Y,d_Y)\) gives levels \((Y,d_{\Lambda,u})\). This is a conventional model for the terminology; the source does not specify whether this model or additional restrictions on the spaces and actions are intended.

For metric-space sequences \((X_n,d_n)\) and \((Y_n,e_n)\), a uniform family of quasi-isometric embeddings consists of maps \(f_n:X_n\to Y_n\) and constants \(A\geq1\), \(B\geq0\), independent of \(n\), such that
\[
 A^{-1}d_n(x,x')-B
 \leq e_n(f_n(x),f_n(x'))
 \leq A d_n(x,x')+B
 \quad(n\in\N,\ x,x'\in X_n).
\]
A uniform family of coarse embeddings consists of maps \(f_n:X_n\to Y_n\) and nondecreasing functions \(\rho_-,\rho_+:[0,\infty)\to[0,\infty)\), independent of \(n\), with \(\lim_{r\to\infty}\rho_-(r)=\infty\), such that
\[
 \rho_-(d_n(x,x'))
 \leq e_n(f_n(x),f_n(x'))
 \leq \rho_+(d_n(x,x'))
 \quad(n\in\N,\ x,x'\in X_n).
\]
Neither definition requires coarse surjectivity.

Under these conventions, the question asks what rigidity conclusion, if any, follows from such embeddings when
\[
 (X_n,d_n)=(X,d_{\Gamma,t_n}),\qquad
 (Y_n,e_n)=(Y,d_{\Lambda,u_n}),\qquad
 t_n\to\infty,\quad u_n\to\infty.
\]
The use of uniform controls and diverging level parameters here explains a possible sequence-level interpretation; neither is explicitly quantified in the supplied question. The relationship between \(t_n\) and \(u_n\), the admissible actions, and the proposed rigidity conclusion remain unresolved. No equivalence of actions, relation between groups, or other specific conclusion can be asserted from this record alone.
\subsection{Short English statement}
The question asks whether embeddings between sequences of warped-cone levels force a rigidity relation between the underlying geometric or dynamical data. It asks about both quasi-isometric and coarse embeddings. The supplied record leaves the necessary hypotheses, the meaning of corresponding levels, and the intended rigidity conclusion unspecified, so this is a statement-only draft requiring human mathematical review.
\subsection{Sources}
\begin{itemize}
\item Local UnsolvedMath registry, supplied complete source record: permanent ID \(30003783\), exact source problem code \texttt{catalogue entry}, title ``Rigidity from Coarse Embeddings of Warped-Cone Levels,'' in the Oberwolfach Reports Open Problems collection. No individual proposer is identified. The supplied proposed year is \(2018\).
\item Supplied report citation: \emph{Mini-Workshop: Superexpanders and Their Coarse Geometry (2019)}, Oberwolfach Reports, DOI \texttt{10.4171/owr/2018/19}, pp.\ 1117--1160. \url{https://doi.org/10.4171/owr/2018/19}. The supplied background attributes the question to workshop open-problem discussions documented in Oberwolfach Reports.
\item Supplied literature source, ``workshop report'': \url{https://ems.press/content/serial-article-files/46739}.
\item Supplied literature source, ``measured-equivalence consequences'': \url{https://arxiv.org/abs/2004.10000}.
\item Supplied literature source, ``dynamical relation from warped cones'': \url{https://doi.org/10.1142/S1793525325500128}.
\item Literature assessment dated exactly \texttt{2026-08-21}: supplied dated curation, not independently verified. The supplied status is ``partially solved.'' The assessment reports stable-product and uniform-measured-equivalence consequences, and examples for quasi-isometric embeddings, under extra hypotheses; it also reports that no general quasi-isometric or coarse-embedding rigidity theorem is established. The background describes this as dated literature triage derived from the supplied Oberwolfach Reports curation and requests independent verification of time-sensitive status and cited claims.
\end{itemize}
\end{document}
